% MODEL-ID: SP-P0001-Q38-M1U-Q4-v1
% STATIC-FLOOR-BYTES: 15660093440
% KV-BYTES-PER-CONTEXT-TOKEN: 65536
% RAW-BUS-BPS: 819200000000
% APPLE-RATED-BPS: 800000000000
% RAW-CTX0-MILLI-TPS: 52311
% APPLE-CTX0-MILLI-TPS: 51085
% RAW-CTX262144-MILLI-TPS: 24945
% APPLE-CTX262144-MILLI-TPS: 24360

\documentclass[11pt,a4paper]{article}
\usepackage[T1]{fontenc}
\usepackage{lmodern}
\usepackage{amsmath,amssymb,mathtools,booktabs,geometry,hyperref,microtype}
\geometry{margin=27mm}
\hypersetup{colorlinks=true,linkcolor=black,urlcolor=blue,citecolor=black}
\newcommand{\bytes}{\operatorname{byte}}
\newcommand{\tps}{\operatorname{token/s}}

\title{A Bandwidth Ceiling for Single-Token Qwen3.8-27B Q4\_K\_M Decoding on M1 Ultra}
\author{Model ID: SP-P0001-Q38-M1U-Q4-v1}
\date{23 August 2026}

\begin{document}
\maketitle

\begin{abstract}
We study batch-one, non-speculative, non-MTP autoregressive decoding of Qwen3.8-27B Q4\_K\_M on M1 Ultra. Rather than identifying the total GGUF file size with DRAM traffic per token, we derive a conservative encoded-byte floor from the architecture and quantization-block formats, then state explicitly the physical premise that actual DRAM traffic is no smaller than this floor. The mathematical argument is a conservation law: if each token requires at least $D$ bytes to cross DRAM and the memory path can carry at most $B$ bytes per second, then throughput must satisfy $rD\le B$. The Lean development measures rate in milli-token/s, eliminating floating-point boundary issues. For short context with \texttt{cacheCarry=0}, the exact discrete ceiling of the abstract model is $51.085$ token/s under Apple's published $800\times10^9$ byte/s bandwidth and $52.311$ token/s under a more permissive $819.2\times10^9$ byte/s raw line-rate scenario. At context length $262144$, the corresponding ceilings are $24.360$ and $24.945$ token/s. Exactness is model-internal; it is not an empirical attainability claim.
\end{abstract}

\section{Scope, quantities, and units}

We consider one ordinary target-model decode step producing one target token. The principal instance fixes batch size one, excludes speculative decoding and MTP, stores the growing full-attention KV cache in FP16, and takes \texttt{cacheCarry=0}. Thus no cross-token on-chip residency is silently deducted from the DRAM floor.

Let $B\in\mathbb N$ denote a memory-bandwidth cap in byte/s, $r\in\mathbb N$ a rate in milli-token/s, $A\in\mathbb N$ the actual DRAM bytes transferred per generated token, and $D(L,C)$ a modeled lower bound on those bytes at context length $L$ with an allowance $C$ for useful on-chip data retained across decode steps.

With rate measured in thousandths of a token per second, bandwidth feasibility is the integer relation
\begin{equation}
  rA\le1000B.
  \label{eq:feasible}
\end{equation}
No floating-point approximation is involved in this definition.

\section{Evidence layer versus theorem layer}

The M1 Ultra bandwidth figure, Qwen3.8-27B architectural dimensions, GGML block encodings, and the tensor-class assumptions of the selected GGUF artifact are external facts. Their provenance is recorded in \texttt{evidence.md}. Lean proves deductions from mathematical premises; it does not authenticate vendor pages or model artifacts.

The physical bridge is therefore stated explicitly:
\begin{equation}
  D(L,C)\le A.
  \label{eq:bridge}
\end{equation}
Every throughput theorem below is conditional on equations~\eqref{eq:feasible} and~\eqref{eq:bridge}.

This is also why the proof does not use ``GGUF file size divided by bandwidth.'' File size is not the same quantity as DRAM traffic for one decode step: files include metadata and alignment, while cache residency may alter memory traffic. We instead construct a smaller tensor-by-tensor representation floor. Omitting additional real traffic only makes the final ceiling more permissive.

\section{Static encoded-byte floor}

For the Qwen3.8-27B text model we use
\begin{equation}
  H=5120,\qquad I=17408,\qquad N_{\mathrm{layer}}=64,
\end{equation}
with 48 Gated DeltaNet layers, 16 full-attention layers, and vocabulary size $V=248320$.

\subsection{Feed-forward networks}
Each SwiGLU FFN contains gate, up, and down matrices. Hence
\begin{equation}
\begin{aligned}
N_{\mathrm{FFN}}
&=64\times3\times5120\times17408\\
&=17\,112\,760\,320.
\end{aligned}
\label{eq:ffn}
\end{equation}

\subsection{Gated DeltaNet projections}
There are 16 Q/K heads and 48 V heads of dimension 128, so
\begin{equation}
 d_{QK}=2048,\qquad d_V=6144.
\end{equation}
The floor includes, per GDN layer,
\[
W_{qkv}:5120\to(2d_{QK}+d_V),\quad
W_z:5120\to d_V,\quad
W_o:d_V\to5120.
\]
Across 48 layers this contributes
\begin{equation}
  N_{\mathrm{GDN,Q4}}=5\,536\,481\,280.
  \label{eq:gdn}
\end{equation}
The selected recipe counts \texttt{in\_proj\_a} and \texttt{in\_proj\_b} as F32, giving
\begin{equation}
  N_{\mathrm{SSM,F32}}=48\times2\times5120\times48
  =23\,592\,960.
  \label{eq:ssm}
\end{equation}

\subsection{Full attention}
Full attention has 24 Q heads, 4 KV heads, and head dimension 256. Because the Q projection also carries the output gate, its output width is counted as $2\times24\times256$. The Q/K/V/O matrices over 16 layers contain
\begin{equation}
  N_{\mathrm{FA,Q8}}=1\,677\,721\,600
  \label{eq:fa}
\end{equation}
weights.

\subsection{Output projection}
The input embedding and language-model output head are untied. A decode step does not stream the entire input embedding table, whereas producing all logits requires the output matrix. We therefore retain
\begin{equation}
  N_{\mathrm{out,Q6}}=248320\times5120
  =1\,271\,398\,400.
  \label{eq:out}
\end{equation}

\subsection{Block-encoded bytes}
We use
\[
\mathrm{Q4_K}:256\to144\ \bytes,
\quad
\mathrm{Q6_K}:256\to210\ \bytes,
\]
\[
\mathrm{Q8_0}:32\to34\ \bytes,
\quad
\mathrm{F32}:1\to4\ \bytes.
\]
All relevant tensor counts are block-aligned. Define
\begin{equation}
\begin{aligned}
D_{\mathrm{static}}
={}&144\frac{N_{\mathrm{FFN}}+N_{\mathrm{GDN,Q4}}}{256}\\
&+34\frac{N_{\mathrm{FA,Q8}}}{32}
+4N_{\mathrm{SSM,F32}}
+210\frac{N_{\mathrm{out,Q6}}}{256}.
\end{aligned}
\label{eq:static}
\end{equation}
Substituting equations~\eqref{eq:ffn}--\eqref{eq:out} gives
\begin{equation}
  \boxed{D_{\mathrm{static}}=15\,660\,093\,440\ \bytes/token}.
  \label{eq:staticvalue}
\end{equation}

Equation~\eqref{eq:staticvalue} is deliberately a lower bound at the representation level. RMSNorm vectors, convolution weights, small state and bias tensors, and kernel-intermediate memory traffic are omitted. Such omissions reduce $D$ and therefore increase the permitted throughput ceiling, which is the conservative direction for an impossibility result.

\section{Context-dependent KV term}

Only the 16 full-attention layers carry a KV cache that grows linearly with context. One historical token contributes
\begin{equation}
\begin{aligned}
D_{\mathrm{KV/token}}
&=16\times4\times256\times2\times2\\
&=65\,536\ \bytes/\text{context-token}.
\end{aligned}
\label{eq:kv}
\end{equation}
Thus
\begin{equation}
  R(L)=15\,660\,093\,440+65\,536L.
  \label{eq:R}
\end{equation}
If $C$ bytes can remain usefully resident on chip between consecutive decode steps, define
\begin{equation}
  D(L,C)=\max\{0,R(L)-C\}.
  \label{eq:D}
\end{equation}
The principal result uses $C=0$. This is not a claim that caches do not exist; it is a refusal to equate an unverified cache-capacity number with useful cross-token residency.

\section{General bandwidth theorem}

\paragraph{Theorem 1.}
Let $B,r,A,D,T\in\mathbb N$. If
\begin{equation}
 D\le A,\qquad rA\le1000B,
 \qquad1000B<TD,
 \label{eq:premises}
\end{equation}
then $r<T$.

\paragraph{Proof.}
Assume instead that $r\ge T$. Since $A\ge D$, monotonicity of multiplication on natural numbers gives $TD\le rA$. Bandwidth feasibility gives $rA\le1000B$, hence $TD\le1000B$. But equation~\eqref{eq:premises} also states $1000B<TD$, a contradiction. Therefore $r<T$.\hfill$\square$

The theorem itself is independent of neural-network details. Architecture is used only to establish a defensible $D$; the final implication is ordinary conservation arithmetic.

\section{Exact discrete ceiling}

If an integer $q$ satisfies
\begin{equation}
  qD\le1000B<(q+1)D,
  \label{eq:exact}
\end{equation}
then in the favorable floor-traffic model $A=D$, the rate $q$ is feasible while $q+1$ is impossible. Hence $q$ is the exact maximum at milli-token/s granularity.

This model-internal feasibility is not a hardware-attainability claim. A real implementation may have $A>D$, fail to sustain the bandwidth cap, or incur additional scheduling and kernel overheads.

\section{Numerical certificates}

\subsection{Apple's 800 GB/s specification}
Set
\[
B_{\mathrm{Apple}}=800\,000\,000\,000\ \bytes/s,
\qquad L=C=0.
\]
Then $D=15\,660\,093\,440$ and
\begin{equation}
51\,085D\le1000B_{\mathrm{Apple}}<51\,086D.
\end{equation}
Therefore
\begin{equation}
  \boxed{r_{\max}=51.085\ \tps}.
\end{equation}

\subsection{The 819.2 GB/s raw line-rate scenario}
Set
\[
B_{\mathrm{raw}}=819\,200\,000\,000\ \bytes/s.
\]
At short context,
\begin{equation}
52\,311D\le1000B_{\mathrm{raw}}<52\,312D,
\end{equation}
so
\begin{equation}
  \boxed{r_{\max}=52.311\ \tps}.
\end{equation}

\subsection{Context length $L=262144$}
Equation~\eqref{eq:R} gives
\begin{equation}
R(262144)
=15\,660\,093\,440+65\,536\times262\,144
=32\,839\,962\,624\ \bytes/token.
\end{equation}
The adjacent-integer certificates yield
\begin{equation}
\boxed{
\begin{aligned}
800\text{ GB/s}:&\quad24.360\ \tps,\\
819.2\text{ GB/s}:&\quad24.945\ \tps.
\end{aligned}}
\end{equation}

\section{Interpretation boundary}

The numerical ceilings are conditional statements, not unconditional laws of the physical M1 Ultra. The 800 GB/s quantity is Apple's product specification. The 819.2 GB/s quantity is a more permissive raw line-rate scenario. Equation~\eqref{eq:staticvalue} is a lower bound derived from frozen tensor classes and block encodings. Equation~\eqref{eq:bridge} is the physical premise connecting real execution to that representation floor.

If future microarchitectural evidence establishes stable cross-token residency for part of the relevant representation, the correct model change is to increase $C$ and re-instantiate the theorem. Speculative decoding and MTP are likewise outside scope because one target verification may correspond to multiple accepted output tokens, changing the measured output-token rate.

\section{Formal verification}

The machine-checkable development is \texttt{Proof.lean}. The intended commands are
\begin{verbatim}
lake build
lake env leanchecker --fresh Problems.Qwen38M1UltraQ4.Proof
lake env lean Problems/Qwen38M1UltraQ4/Proof.lean
\end{verbatim}
CI checks the reported axiom set against a subset of Lean's standard $\{\texttt{propext},\texttt{Classical.choice},\texttt{Quot.sound}\}$ set and rejects \texttt{sorryAx} and \texttt{Lean.trustCompiler}.

The claim certified by Lean is therefore best written as
\[
\boxed{
\text{external evidence and physical premises}
\Longrightarrow
\text{the stated discrete throughput ceilings}}
\]
rather than treating the evidence itself as if it had been proved in the kernel.

\end{document}
